\documentclass[11pt]{article}
\usepackage[margin=1in]{geometry}
\usepackage[T1]{fontenc}
\usepackage{lmodern,amsmath,amssymb,amsthm,mathtools,hyperref}
\setlength{\emergencystretch}{3em}
\newtheorem{theorem}{Theorem}[section]
\newtheorem{lemma}[theorem]{Lemma}
\newtheorem{corollary}[theorem]{Corollary}
\newtheorem{proposition}[theorem]{Proposition}
\newcommand{\E}{\mathbb E}
\newcommand{\tr}{\operatorname{tr}}
\title{Support rigidity, optimal atom exchange,\\and modular Hadamard-extension counts}
\author{Joshua Mark\thanks{ORCID: \href{https://orcid.org/0009-0003-5905-2382}{0009-0003-5905-2382}. AI tools assisted with derivation, independent review, computation and exposition.}}
\date{20 November 2026 --- research draft; priority not established}
\begin{document}\maketitle

\begin{abstract}
We continue the search for concentrated pairwise-independent laws.
An elementary binary-rank theorem forces almost complete support
replacement for low-rank candidates, even along known Hadamard orders
where moment errors vanish and nearby feasible laws exist. We give
sharp support-dependent Boolean degree criteria, including the
punctured extraction target. In the four-bit moment polytope, an
explicit circuit jump reaches an optimum and its exact barrier is
optimal among all continuous feasible paths. A modular lifting
theorem converts residue counts into exact sign-extension counts by
a column-pair correction in a specified eight-divisible order class.
These are conditional tools, finite-model constructions and structural
bounds. They do not construct the required distribution at an unknown
Hadamard order or prove the full conjecture. Priority remains unverified.
\end{abstract}
The standing extraction target is an exactly unbiased pairwise-independent
law on $N-4$ bits with collision at least $1/N-1/(1024N^4)$, for
$N\ge8$ divisible by four. The proof is preserved in
\path{../round6/collision_stability.tex}. Its exact entropy and
minimal-support prerequisites are in
\path{../round5/collision_entropy.tex}. We do not assume these inputs
exist in every dimension.
\section{Small moment errors can require almost complete support replacement}
\label{sec:rank-rounding}
This section connects the distribution target with elementary binary
matrix rigidity. We distinguish binary conversion $(1-H)/2$ from
reducing the sign matrix itself modulo two, which would lose all
sign information.

Let $N\ge12$, $N\equiv4\pmod8$, and put $d=N-4$. The preceding
collision theorem says that an exactly unbiased pairwise-independent
law on at most $N$ sign vectors is uniform on exactly $N$ vectors.
Thus every minimal-support target has the following binary rank.

\begin{lemma}[Binary rank of a minimal-support target]
\label{lem:rank-target}
For such a target, let $Y$ be the $N\times d$ binary matrix with
$Y_{ij}=(1-x_{ij})/2$. Then
\[
 Y^TY=I_d+J_d\quad\text{over }\mathbb F_2,
 \qquad \operatorname{rank}_{\mathbb F_2}Y=d.
\]
\end{lemma}
\begin{proof}
Every column has $N/2$ negative signs, an even integer. Every distinct
pair of columns has $N/4$ simultaneous negative signs, an odd integer.
These give diagonal zero and off-diagonal one. Since $d$ is even,
$(I+J)^2=I$ over $\mathbb F_2$, so the Gram matrix is invertible.
\end{proof}

\begin{theorem}[Support-overlap rigidity]\label{thm:rank-overlap}
If a binary linear code $\mathcal C\subseteq\mathbb F_2^d$ has
dimension at most $r$, every minimal-support target intersects it in
at most $r+4$ atoms. Consequently a uniform $N$-point candidate
supported in $\mathcal C$ has total-variation distance at least
\[
 1-\frac{r+4}{N}
\]
from every exactly pairwise-independent law on at most $N$ points.
At least $N-r-4$ atoms must be replaced; changing their weights alone
cannot avoid this conclusion.
\end{theorem}
\begin{proof}
Let $k$ target rows belong to $\mathcal C$. Their binary span has
dimension at most $r$, and the remaining $N-k$ rows add at most
$N-k$ dimensions. Lemma~\ref{lem:rank-target} gives
$N-4\le r+N-k$, hence $k\le r+4$.
Both laws in the distance assertion are uniform on $N$ points;
their total-variation distance is one minus their support overlap
divided by $N$. The earlier support theorem also rules out a smaller
positive support or nonuniform exact reweighting.
\end{proof}

\begin{proposition}[An explicit near-moment candidate]
\label{prop:rank-walsh-candidate}
Let $L=2^r$, $d=N-4\ge L/2$, and $N\le L$. Observe the distinct
characters with binary masks $1,\ldots,d$ on $\mathbb F_2^r$.
Choose any $N$ seed points and give them equal weights. The resulting
law $P$ has $N$ distinct atoms,
\[
 C(P)=1/N,\qquad
 \max_{1\le|S|\le2}|\mathbb E_P\chi_S|\le\frac{L-N}{N}.
\]
Nevertheless its distance from every minimal-support target is at
least $1-(r+4)/N$.
It is simultaneously close to an exactly feasible law: if $U$ is
the full uniform Walsh law, then
\[
 d_{\rm TV}(P,U)=1-N/L,\qquad C(U)=1/L.
\]
\end{proposition}
\begin{proof}
The observed masks include a binary basis because $d\ge L/2$, so
the seed-to-atom map is injective. Under the full seed law every
mean and distinct pair correlation is zero. Removing $L-N$ seed
points changes each unnormalized character sum by at most $L-N$,
giving the stated bound. The atoms lie in a binary linear code of
dimension $r$, so apply Theorem~\ref{thm:rank-overlap}. Direct counting
gives the distance to $U$.
\end{proof}
In particular, $N=2^r-4$ gives moment error at most $4/N$, while the
fraction of atoms that must be replaced tends to one. This statement
does not require assuming existence at every order $2^r-4$.

\begin{theorem}[A nonvacuous asymptotic obstruction to dimension-free rounding]
\label{thm:rank-asymptotic}
There is an infinite sequence of classical known Hadamard orders $N$
and laws $P_N$ on $N-4$ bits such that
\begin{enumerate}
\item $C(P_N)=1/N$ and every first and second moment error tends to zero;
\item $P_N$ is at total-variation distance tending to zero from an
exactly unbiased pairwise-independent law;
\item its distance from every exact law on at most $N$ atoms tends to one.
\end{enumerate}
The exactly feasible laws in item two also have
$N C(U_N)\longrightarrow1$. Thus even vanishing moment error and
asymptotically optimal \emph{relative} collision do not give a
dimension-free total-variation rounding theorem to extremizers.
\end{theorem}
\begin{proof}
Paley's classical theorem supplies a Hadamard matrix of order
$N=q+1$ for every prime power $q\equiv3\pmod4$
(\emph{On Orthogonal Matrices}, 1933,
\url{https://doi.org/10.1002/sapm1933121311}).
Take $q=3^{2t+1}$, so $N\equiv4\pmod8$.
Irrationality of $\log_2 9$ makes the fractional parts of
$(2t+1)\log_2 3$ dense in the unit interval. Select an increasing
subsequence approaching one from below and let
$L=2^{\lceil\log_2 q\rceil}$. Then $q/L\to1$. Since $q\equiv3\pmod8$,
we have $q\le L-5$ for large $L$, so $L\ge N+4$ and $N/L\to1$.
Eventually $d=N-4\ge L/2$. Apply the preceding proposition.
Here $r=\log_2 L=O(\log N)$, so $(r+4)/N\to0$.
The target class is nonempty because these are Paley orders.

The elementary density fact used here follows from pigeonholing the
multiples of an irrational angle to obtain arbitrarily small nonzero
rotation steps. Their positive multiples form arbitrarily fine nets
of the circle, and the same holds after translating or discarding any
finite prefix of the orbit. No claim about primes near powers of two
or unproved prime-pattern conjecture is used.
\end{proof}

This is a restriction on concentration-preserving rounding and on
repairs changing few atoms. It is not a failure of moment correction:
the full Walsh law is explicitly nearby and feasible. Nor does it
contradict the earlier dimension-dependent extraction theorem, whose
collision window is much smaller. The elementary rank methods here
give a quantitative diagnostic; their bibliographic novelty is not
asserted.
\section{Fixed-support degree criteria and punctured-observable obstructions}
\label{sec:nonlinear-fixed-support}
We use Boolean algebraic degree over $\mathbb F_2$, not real Walsh degree.
Let $S\subset\mathbb F_2^r$ have size $N$, let $T$ be its nonempty
complement, and suppose $N\equiv4\pmod8$, $r\ge4$. The seed law is
uniform on $S$. Put $t=\operatorname{affdim}(T)$ and $a=r-t$.
The low-degree model considered first is degree at most $r-2$.
The code background is standard Reed--Muller duality; see Carlet's
\emph{Boolean Functions for Cryptography and Coding Theory},
\url{https://www.math.univ-paris13.fr/~carlet/book-fcts-Bool-vect-crypt-codes.pdf}.
The duality and all required rank facts are proved below. No novelty
or new Hadamard order is asserted.

\begin{lemma}[Even restriction code and its radical]
\label{lem:nonlinear-radical}
Let $C=\mathrm{RM}(r-2,r)|_S$ be the binary restriction code, with
$C_0$ its even-weight subcode. Then
\[
 \dim C=N-a,\quad \dim C_0=N-a-1,\quad
 \dim\operatorname{rad}(C_0)=a+1,
\]
where the radical uses the binary dot product. Consequently its bilinear
rank is $N-2a-2$.
\end{lemma}
\begin{proof}
Monomials of combined degree less than $r$ have even full-cube sum;
dimension counting gives $\mathrm{RM}(r-2,r)^\perp=\mathrm{RM}(1,r)$.
Thus $C^\perp$ consists of the restrictions of affine polynomials
vanishing on $T$, a space of dimension $r-t=a$. Every such restriction
lies in $C$ and has even weight: its nonzero full-cube weight is $2^{r-1}$.
Also $\mathbf1\in C$ has even weight $N$.
The vector $\mathbf1$ is not in $C^\perp$, because that would make
$\mathbf1_S$ affine on the full cube, impossible at weight
$N\equiv4\pmod8$ with $r\ge4$. Hence the parity functional on $C$ is
nonzero, giving $\dim C_0=N-a-1$. Finally
$C_0^\perp=C^\perp+\langle\mathbf1\rangle\subseteq C_0$;
therefore the radical is this entire $(a+1)$-dimensional space.
\end{proof}

\begin{theorem}[Exact fixed-support criterion for a full basis]
\label{thm:nonlinear-fixed-support}
Suppose an order-$N$ Hadamard matrix exists. The support $S$ admits
$N-1$ unbiased pairwise-independent sign observables of degree at most
$r-2$ if and only if $T$ affinely spans $\mathbb F_2^r$.
If $T$ does span, every prescribed row labeling of any such Hadamard
matrix admits observables with this degree bound.
\end{theorem}
\begin{proof}
Let $B$ be the binary truth columns of the observables on $S$.
Balancedness gives weight $N/2$, even; pairwise orthogonality gives
intersection size $N/4$, odd. Hence $B^TB=I+J$ over $\mathbb F_2$.
For $N-1$ columns its rank is $N-2$. All columns lie in $C_0$, so
Lemma~\ref{lem:nonlinear-radical} forces
$N-2\le N-2a-2$, or $a=0$.
Conversely, $a=0$ means $C^\perp=0$, so restriction of
$\mathrm{RM}(r-2,r)$ to $S$ is surjective. Extend each labeled binary
Hadamard column separately.
\end{proof}

\paragraph{An explicit support-switching operation.}
This surjectivity can be implemented as a small syndrome repair, rather
than a new interpolation search. Relabel the existing rows arbitrarily
onto $S$, set provisional outside values to zero, and choose an affine
basis $z_0,\ldots,z_r$ inside $T$. Compute the $r+1$ binary parity defects
of the provisional truth table against $1,z_1,\ldots,z_r$ (coordinate
functions, with the basis points distinguished by their indices).
Flip values only at the chosen basis points to cancel those defects.
The matrix with columns $(1,z_i)$ is invertible, so there is a unique
repair on those points. The repaired function is orthogonal to all
affine functions and therefore has degree at most $r-2$.
No value on $S$ changes, so all real moments are preserved exactly.
For degree $r-1$, a single omitted point suffices: repair only total
truth-table parity. This operation works for any proper support and
permits arbitrary row relabeling.

\begin{corollary}[Degree saturation near a power of two]
\label{cor:nonlinear-degree-saturation}
Assume $2^{r-1}<N<2^r$, $N\equiv4\pmod8$, and $H_N$ exists.
There is some minimal-seed encoding of degree at most $r-2$ if and
only if $2^r-N\ge r+1$. If $2^r-N\le r$, the smallest possible maximum
degree is exactly $r-1$. In particular every existing
$H_{2^r-4}$, $r\ge4$, requires degree exactly $r-1$ on a minimal seed.
\end{corollary}
\begin{proof}
A set of at most $r$ omitted labels cannot affinely span an $r$-space;
a set of at least $r+1$ can be chosen to contain an affine basis.
Use the theorem, and the one-point parity repair for the upper bound.
\end{proof}

\subsection{The $N-4$-bit extraction input still requires growing degree}
\begin{theorem}[Punctured-observable bounds, without completion]
\label{thm:nonlinear-punctured}
Let $N\equiv4\pmod8$ and let $N$ equal-weight seed labels carry $m=N-4$
unbiased pairwise-independent sign bits, all of Boolean degree at most
$d$. Then
\[
 \sum_{j=0}^d\binom rj\ge N-3.
\]
At minimal $r=\lceil\log_2N\rceil$, this forces
$d\ge\lceil r/2\rceil$. For degree at most $r-2$, the sharper necessary
condition on the omitted labels is $\operatorname{affdim}(T)\ge r-1$.
For $N=2^r-4$ and $r\ge5$, degree $r-2$ is impossible for this punctured
input, not only for a full Hadamard basis.
\end{theorem}
\begin{proof}
The same binary Gram is $I+J_m$, now nonsingular because $m$ is even.
Thus the $m$ truth columns are independent. The constant vector is
orthogonal to them, is nonzero, and cannot belong to their nondegenerate
span. The Boolean restriction space must have dimension at least
$m+1=N-3$. Since $N-3>2^{r-1}$ at minimal seed length, binomial symmetry
proves the growing-degree bound.
For $d=r-2$, the nonsingular Gram embeds an $m$-dimensional nondegenerate
space into $C_0$. Its bilinear rank must therefore be at least $m$,
so $N-4\le N-2a-2$ implies $a\le1$. Four omitted points have affine
dimension at most three; when $r\ge5$, this violates $t\ge r-1$.
No completion theorem or existing full Hadamard matrix was used.
\end{proof}

\subsection{A sharp four-bit quadratic classification}
\begin{proposition}[Exact finite obstruction]\label{prop:nonlinear-twelve-classification}
For a uniform law on any twelve labels in $\mathbb F_2^4$, the largest
number of unbiased pairwise-independent quadratic sign observables is:
\[
 \begin{cases}
 5,&\text{if the four omitted labels form an affine plane},\\
 7,&\text{if they are affinely independent}.
 \end{cases}
\]
In particular the order-twelve $N-4=8$-bit extraction input requires
at least one cubic observable, for every labeling.
\end{proposition}
\begin{proof}[Exact finite certificate]
Affine input changes preserve degree and real moments. Four distinct
points in $\mathbb F_2^4$ have precisely two affine equivalence types:
representatives $\{0,1,2,3\}$ and $\{0,1,2,4\}$.
For each representative, enumerate the binary code generated by all
degree-at-most-two monomial restrictions to the twelve retained labels.
It has respectively $1024$ and $2048$ words. Balanced sign observables
are exactly weight-six words; identify a word with its complement by
requiring its first coordinate zero. This leaves respectively $126$
and $238$ vertices. Join two words when their xor has weight six.
Cliques are exactly sets of unbiased pairwise-independent observables.
The accompanying original standard-library checker exhaustively excludes
cliques of sizes six and eight, respectively, and verifies cliques of
sizes five and seven. Sorted-support integer-word witnesses are
\[
 \begin{array}{c|l}
 \{0,1,2,3\}&246,826,1372,1708,3178\\
 \{0,1,2,4\}&246,826,1372,1708,3178,3504,3606.
 \end{array}
\]
All enumeration and recursive candidate intersections are exact bitset
operations; the checker also recovers explicit quadratic polynomials.
This is an exhaustive finite computational certificate, not a numerical
optimization or a Lean formalization.
\end{proof}
Combining the proposition with Theorem~\ref{thm:nonlinear-punctured}
shows that the punctured input at every order $2^r-4$, $r\ge4$, requires
maximum degree at least $r-1$ in this equal-weight minimal-seed model.
Whenever a full Hadamard matrix at that order exists, one-point parity
repair gives matching degree $r-1$ examples. This obstruction does not
apply automatically to nonuniform measures, larger seeds, or approximate
moments; those are distinct possible search directions.

\section{An optimal atom exchange and its exact global path barrier}

For four unbiased pairwise-independent sign bits, the free Walsh
coefficients are $a\in\mathbb R^5$, comprising the four triple
characters and the quadruple character. The feasible polytope is
\[
 P=\{a:1+v\cdot a\ge0\text{ for all }v\in\{\pm1\}^5,
 \ \prod_i v_i=1\}.
\]
The map from bit strings to these sixteen $v$ is a bijection.
The previously established vertex description is
$\{\pm e_i\}_{i=1}^5\cup\{s/3:\prod_i s_i=1\}$.
Its collision probability is $C(a)=(1+\|a\|^2)/16$.
The ten first vertices give uniform eight-atom laws with collision
$1/8$; the sixteen latter vertices have one atom of weight $1/6$,
ten of weight $1/12$, and collision $7/72$.

\begin{theorem}[A finite exact escape rule]
At any weighted vertex $a=s/3$, choose any $i\in\{1,\ldots,5\}$
and replace the law by the uniform vertex $b=s_i e_i$.
This is an edge move of the feasible moment polytope, preserves all
normalization, positivity and pairwise constraints, and changes
twelve atom weights. It increases collision by exactly $1/36$.
Along its feasible interpolation $a(t)=(1-t)a+tb$,
\[
 C(a(t))=\frac7{72}-\frac{t}{36}+\frac{t^2}{18}.
\]
It initially decreases collision, returns to the initial value at
$t=1/2$, and reaches its minimum $3/32$ at $t=1/4$.
\end{theorem}
\begin{proof}
The weighted vertex has five active inequalities, with normals
$s$ multiplied coordinatewise by each sign vector having four minus
signs. At $b$, all except the normal whose plus position is $i$
remain active. These four normals are independent, so their common
face has dimension at most one. The segment between the two distinct
vertices lies in that face, hence is an edge. Its midpoint is a
strict convex combination of two laws; no nonnegativity is lost.
The weighted support has eleven atoms, the uniform support eight,
and their intersection seven, leaving a twelve-atom union.
All twelve weights change: the nonzero differences have magnitudes
$1/24$, $1/12$ or $1/8$. Expanding
$\|a(t)\|^2=5/9-4t/9+8t^2/9$ proves the formula.
The edge is a circuit in the normalization/pairwise moment kernel:
its support-restricted kernel has dimension one, and every entry of
its direction on those twelve atoms is nonzero.
\end{proof}

\begin{theorem}[The unavoidable barrier is exactly $1/288$]
Every continuous feasible path from a weighted vertex to any uniform
vertex passes through a law of collision at most $3/32$.
The preceding edge path attains this upper bound on its minimum, so
it is optimal in the maximin sense. In particular no continuously
collision-nondecreasing path escapes a weighted vertex to an optimum.
\end{theorem}
\begin{proof}
We first show that every $a\in P$ with $|a_i|=1/2$ satisfies
$\|a\|^2\le1/2$. Permute coordinates to take $i=1$.
If $a_1=1/2$, write $y=(a_2,\ldots,a_5)$. The constraints with
$v_1=-1$ use the eight four-dimensional odd-parity sign vectors.
They consist of $\pm h_j$, where the four rows $h_j$ of $J_4-2I_4$
are mutually orthogonal with squared norm four. Thus
$|h_j\cdot y|\le1/2$, and
\[
 \|y\|^2=\frac14\sum_{j=1}^4(h_j\cdot y)^2\le1/4.
\]
If $a_1=-1/2$, use $v_1=+1$ instead; the even-parity four-sign
vectors likewise consist of the signed rows of a Hadamard matrix
of order four. The identical bound follows. More generally the
argument shows $\|a\|^2\le a_i^2+(1-|a_i|)^2$.

At a weighted vertex all five coordinate magnitudes are $1/3$;
at a uniform vertex one is one. The intermediate value theorem
forces some coordinate to cross magnitude $1/2$ along every
continuous connecting path. At that point collision is at most
$(1+1/2)/16=3/32$. The displayed edge path has minimum exactly that
value. Since $7/72-3/32=1/288$, the claimed barrier follows.
\end{proof}

This is an actual nonlocal twelve-atom correction rule on the fully
solved four-bit testbed. It demonstrates that strict local collision
maxima can require a quantified loss before an exact feasible path
reaches a better distribution. It does not assert that every higher
dimensional moment polytope has such a bounded circuit or favorable
endpoint. The four-dimensional Hadamard frame used to certify the
barrier is supplied explicitly; no unknown-order construction is
assumed or obtained.
\section{Exact lifting of modular sign extensions with a column-pair correction}
\label{sec:modular-sign-lifting}

Finite-field residue constraints can retain information that an
all-odd integer kernel test misses. They also admit exact generating
polynomial and Fourier counts. Counting through code enumerators and
character sums is classical; a foundational primary source is
MacWilliams, \emph{A Theorem on the Distribution of Weights in a
Systematic Code} (1963),
\url{https://doi.org/10.1002/j.1538-7305.1963.tb04003.x}.
Here we derive a specific lifting theorem: in a range of dimensions,
every spurious modular sign extension is a two-coordinate perturbation
of an existing row. A simple column statistic subtracts all such
spurious solutions exactly. No novelty claim is made pending a
targeted prior-art review.

Let $P$ have $r$ orthogonal sign rows of length $N$. A modular sign
extension modulo $q$ is a vector $v\in\{\pm1\}^N$ satisfying
$Pv\equiv0\pmod q$; an exact extension satisfies $Pv=0$ over the
integers. Count both kinds up to negation. Write their counts as
$M_q(P)$ and $E(P)$, and put
\[
 A(P)=\#\left\{\{a,b\}:a<b,\quad
 \left|\sum_{i=1}^r p_{ia}p_{ib}\right|=r-2\right\}.
\]
This statistic is computable in $O(rN^2)$ integer operations.

\begin{theorem}[Exact modular count with a column-pair correction]
Suppose $q\ge3$ is odd,
\[
 N=4(q+1),\qquad r>\frac{4(q+1)^2}{q^2}.
\]
Then
\[
 \boxed{E(P)=M_q(P)-A(P).}
\]
More precisely, every modular sign extension which is not exact differs
in exactly two coordinates from the sign or negative of a unique
existing row. Such nonexact classes correspond bijectively to the
column pairs counted by $A(P)$.
\end{theorem}
\begin{proof}
Let $d_i=p_i\cdot v$. Any two existing rows have Hamming distance
$N/2=2(q+1)$, which is even. Hence their sign products have the same
parity, and all $d_i$ are congruent modulo $4$: either all are zero
modulo $4$, or all are $2$ modulo $4$.
Each $d_i$ is even and divisible by odd $q$. In the second case,
$|d_i|=2q$, because the next possible magnitude $6q$ exceeds
$N=4q+4$. Bessel's inequality would then give
\[
 4q^2r=\sum_i d_i^2\le N\|v\|^2=N^2,
\]
contradicting the row-count hypothesis.
Thus every $d_i$ is divisible by $4q$, and its only possibilities
are $0,4q,-4q$. Two nonzero values would contribute
$32q^2>N^2$, again contradicting Bessel. If all vanish, the extension
is exact. Otherwise there is a unique index $i$ with
$|p_i\cdot v|=4q=N-4$. After negating $v$ if necessary, it differs
from $p_i$ in exactly two coordinates $a,b$.

For every $j\ne i$, orthogonality to $v$ and $p_i$ gives
$p_{ja}p_{ia}+p_{jb}p_{ib}=0$. Consequently the column products
$p_{ja}p_{jb}$ have the same sign for all $j\ne i$ and its opposite
sign for row $i$. Their sum has absolute value $r-2$.

Conversely, each column pair with this property has a unique exceptional
row $i$, since the hypotheses imply $r>2$. Flip those two coordinates
of $p_i$. The resulting vector has inner product $N-4=4q$ with
$p_i$ and zero with every other row, so is a nonexact modular extension.
No resulting class can come from two different rows: it would then
lie within Hamming distance two of signs of two orthogonal rows,
whose mutual distance is $N/2\ge8$, whereas the triangle inequality
would bound it by four. The flipped pair is unique once the row is
fixed. This proves the bijection and count formula.
\end{proof}

In particular, if all distinct columns have absolute correlation
strictly below $r-2$, then modular sign extensions are automatically
exact. This is a sufficient \emph{lifting} theorem, not an assertion
that modular solutions always exist.

\subsection{An exact algebraic certificate and positivity criterion}
When $q$ is prime, let $c_j\in\mathbb F_q^r$ be column $j$ of $P$.
In the integral group algebra of $\mathbb F_q^r$, form
\[
 \Phi_P=\prod_{j=1}^N(X^{c_j}+X^{-c_j}),\qquad
 C_q(P)=[X^0]\Phi_P.
\]
Expanding the product counts sign choices with zero residue, so
$C_q(P)=2M_q(P)$. Under the theorem's hypotheses,
\[
 \boxed{E(P)=\tfrac12 C_q(P)-A(P).}
\]
Thus the exact computable inequality $C_q(P)>2A(P)$ suffices for a
sign extension; equality certifies that none exists. Character
orthogonality gives the equivalent Fourier expression
\[
 C_q(P)=\frac1{q^r}\sum_{a\in\mathbb F_q^r}
          \prod_{j=1}^N2\cos\left(\frac{2\pi\,a\cdot c_j}{q}\right).
\]
The group-algebra expression permits integer computation without
floating-point trigonometry. Alternatively modular row reduction of
rank $s$ enumerates at most $2^{N-s-1}$ sign assignments to information
coordinates. The structural correction avoids an additional integer
orthogonality test for every surviving modular candidate.

For a nonvacuous correction example, take rows $0,1,2,3,4,5,6,8$ of
the order-$16$ Sylvester matrix. With $q=3$ one obtains
\[
 C_3(P)=88,\qquad M_3(P)=44,\qquad A(P)=8,\qquad E(P)=36.
\]
The verifier independently computes $C_3(P)$ by integer group-algebra
multiplication and by modular information-coordinate enumeration.

\subsection{The order-twenty-four repair benchmarks}
For $N=24$, use $q=5$. The row-count condition is simply $r\ge6$.
The following exact results concern the two previously recorded
terminal frames and the two single-row deletions of the twelve-row
frame. All counts are antipodal classes.
\[
\begin{array}{c|r|r|r|r|r}
\text{frame}&r&\max_{a<b}|c_a\cdot c_b|&M_3&M_5&E\\\hline
\text{terminal twelve-row}&12&8&72&0&0\\
\text{terminal sixteen-row}&16&4&76&0&0\\
\text{delete row zero}&11&7&72&2&2\\
\text{delete row one}&11&7&182&16&16
\end{array}
\]
In every row the column statistic $A(P)$ is zero, explaining why the
modulo-$5$ candidate sets coincide exactly with the real extension
sets. Modulo $3$ alone instead gives many false positives. In
particular the twelve-row frame passes all-odd integer feasibility
and modulo-$3$ sign feasibility, but fails modulo-$5$ sign feasibility.
This is a concrete stronger residue obstruction, not a denominator
or parity obstruction.

The sixteen-row frame has rank $16$ modulo $5$, so its nonexistence
certificate tests only $2^7=128$ information assignments with elementary
finite-field arithmetic. The twelve-row frame uses $2^{11}=2048$.
The two versus sixteen exact candidates after deletion also recover
the different repair possibilities, after testing their mutual
orthogonality. Candidate existence alone still does not ensure that
enough candidates are compatible for an augmenting exchange.

\subsection{A general moment bound outside the lifting dimension}
There is a weaker sufficient condition without the special relation
$N=4(q+1)$. For even $N$ and odd $q$, let $S_q$ be the modular sign
classes and put
\[
 T_q=|S_q|,\qquad B_q=\sum_{v\in S_q}\|Pv\|^2.
\]
A nonexact modular class has at least one nonzero inner product, an
even multiple of $q$, so its energy is at least $4q^2$. Consequently
\[
 E(P)\ge T_q-\frac{B_q}{4q^2}.
\]
In particular $B_q<4q^2T_q$ certifies an exact extension. This first
moment estimate is often weak; the theorem above succeeds by fully
classifying the possible nonzero energies and their column-pair
geometry. Neither criterion currently proves positivity uniformly
for all admissible orders.

\section{What still needs to be constructed}
The rank example is close to an exactly feasible Walsh law, but that
nearby law has too many equal-weight atoms. It is far from every
minimal-support extremizer. The nonlinear degree criteria identify
which fixed-support representations can work, conditional on
Hadamard existence; they do not supply the required row values at
unknown orders. The exact four-bit exchange solves a calibrated
finite landscape, not a general atom-replacement problem.

The modular theorem offers a concrete extension-count test, but no
uniform positivity theorem for that count has been proved. Its current
arithmetic class is divisible by eight. A full proof still needs a
construction or positivity argument at every admissible order, with
the four-modulo-eight class included. No theorem here assumes away
that missing step. All finite checks are exact arithmetic; they are
not Lean formalization. No Git commit was made.
\end{document}
